\documentclass[11pt]{article}
\usepackage{graphicx} 
\usepackage[a4paper, margin = 1.0in]{geometry}
\usepackage{setspace}
\usepackage{authblk}
\usepackage{indentfirst}
\onehalfspacing


\title{\textbf{Iba't Ibang Statistical Tests Para sa Pananaliksik}}
\author{
    Zion Joseph Sabangan Busto \\
    Instituto ng Estadistika\\
    \texttt{zsbusto@up.edu.ph}
}



\date{November 2026}




\begin{document}
\begin{figure}
    \centering
    \includegraphics[width=0.3\linewidth]{University_of_the_Philippines_seal_2017.png}
\end{figure}



\maketitle
\thispagestyle{empty}
\newpage
\tableofcontents
\pagenumbering{roman}
\newpage
\pagenumbering{arabic}



\section{Introduksyon}

Isa ang estadistika sa pinakamahalagang larangan sa buong mundo. Ito ay ginagamit araw-araw ng mga tao upang gumawa ng mga desisiyon, prediksiyon, at iba pa. Sa panahon ngayon, mahalagang malaman ng mga Pilipino ang iba't ibang konsepto sa estadistika upang sila ay makagawa ng lohikal na pagsusuri. 


\section{Mga Konsepto sa Estadistika}

Sa parte ng papel na ito, ang mananaliksik ay nagpresenta ng ilang statistical methods na simple at  puwedeng magamit kahit na may onting kaalaman lamang sa Sipnayan.

\subsection{Hypothesis Testing}

\subsection{Analysis of Variance}

\subsection{Regression Analysis}

Ang regression analysis ay isang statistical method kung saan ay gumagamit tayo ng equation upang magamit ito para ipaliwanag ang isang bagay. Halimbawa, gusto nating i-analyze ang 

\section{Konklusiyon}

\end{document}
